\section{Appendix: proofs and verification scope}

\subsection{Estimator and uniform upper bound}

The observed category-cell masses specify the inputs to the exact Jackson--factorial construction.

\begin{definitionv}{synth_8}[Observed category-cell masses]
For a covariate cell \(j\) and a treatment-outcome category \(\zeta=(a,y)\in\{0,1\}^2\), define
\[
q_{\zeta j}=P(X=j,A=a,Y=y).
\]
These four masses are the components of the observed cell vector \(q_j=(q_{\zeta j}:\zeta\in\{0,1\}^2)\).
\end{definitionv}

\begin{algorithmv}{def:jf-estimator}[Jackson–factorial estimator $\widehat V_b^{\mathrm{JF}}$]
Given \(O^{(n)}\), \(d\), and \(I\), the estimator \(\widehat V_b^{\mathrm{JF}}\) uses the tuning quantities \(H_0,\kappa,D_0,L,m,K\) defined below.
\begin{enumerate}
\item Set
\[
H_0=4096,\qquad
\kappa=\frac1{512},\qquad
D_0=16,\qquad
L=\log(ed),\qquad
m=\frac n8,\qquad
K=\max\left\{2,\left\lfloor\frac L{512}\right\rfloor\right\},
\qquad
\mathcal Z=\{0,1\}^2.
\]
Index treatment-outcome categories by \(\zeta=(a,y)\in\mathcal Z\).

\item If \(d<D_0\), form the empirical category-cell masses and threshold process:
\[
\widehat q_{\zeta j}
=
\frac1n\sum_{i=1}^n
\mathbf1\{X_i=j,\ (A_i,Y_i)=\zeta\},
\qquad
\widehat q_j=(\widehat q_{\zeta j})_{\zeta\in\mathcal Z},
\qquad
\widehat F(\lambda)=\sum_{j=1}^d f_\lambda(\widehat q_j).
\]
Proceed to the final step.

\item If \(d\geq D_0\) and \(n<d/L\), output
\[
\widehat V_b^{\mathrm{JF}}=\frac12,\qquad b\in I.
\]
In the remaining branch, \(d\geq D_0\) and \(n\geq d/L\), draw \(M\sim\operatorname{Pois}(n/4)\) independently of \(O^{(n)}\). On \(\{M>n\}\), set
\[
\widetilde F^{\mathrm{cap}}(\lambda)=0,\qquad \lambda\in[0,1].
\]
On \(\{M\leq n\}\), uniformly permute the records, retain the first \(M\), and independently mark each retained record pilot or evaluation with probability \(1/2\) each. Define the pilot and evaluation counts by
\[
N'_{\zeta j}
=
\sum_{i=1}^{M}
\mathbf1\{X_i^*=j,\ (A_i^*,Y_i^*)=\zeta,\ i\text{ is marked pilot}\},
\qquad
N_{\zeta j}
=
\sum_{i=1}^{M}
\mathbf1\{X_i^*=j,\ (A_i^*,Y_i^*)=\zeta,\ i\text{ is marked evaluation}\},
\]
where \(O_i^*=(X_i^*,A_i^*,Y_i^*)\) denotes the \(i\)th record in the uniform permutation.

\item On \(\{M\leq n\}\), for every \(j\) and \(\zeta\), set
\[
\delta=\frac Lm,\qquad
c_{\zeta j}=\frac{N'_{\zeta j}}m,\qquad
h_{\zeta j}=H_0\left\{\sqrt{c_{\zeta j}\delta}+\delta\right\},
\]
\[
\ell_{\zeta j}=\max\{0,c_{\zeta j}-h_{\zeta j}\},
\qquad
u_{\zeta j}=c_{\zeta j}+h_{\zeta j},
\qquad
c^0_{\zeta j}=\frac{\ell_{\zeta j}+u_{\zeta j}}2,
\qquad
R_{\zeta j}=\frac{u_{\zeta j}-\ell_{\zeta j}}2,
\qquad
S_j=\sum_{\zeta\in\mathcal Z}R_{\zeta j}.
\]
All radii \(R_{\zeta j}\) are strictly positive. Let \(c_j^0=(c^0_{\zeta j})_{\zeta\in\mathcal Z}\) and \(R_j=(R_{\zeta j})_{\zeta\in\mathcal Z}\). Define the normalized positive order-four Jackson kernel by
\[
J_K(t)=C_K\left\{\frac{\sin(Kt/2)}{\sin(t/2)}\right\}^{4},
\qquad -\pi\leq t\leq\pi,
\qquad
\int_{-\pi}^{\pi}J_K(t)\,dt=1,
\]
where the quotient takes its continuous value at zero and \(C_K>0\) is the unique normalizing constant. Set \(D=2(K-1)\). For each \(j\) and \(\lambda\), define the unique tensor polynomial \(P_{j,\lambda}\), of coordinatewise degree at most \(D\), through
\[
P_{j,\lambda}(c_j^0+R_j\odot\cos\vartheta)
=
\int_{[-\pi,\pi]^4}
f_\lambda\!\left(c_j^0+R_j\odot\cos(\vartheta+t)\right)
\prod_{\zeta\in\mathcal Z}J_K(t_\zeta)\,dt,
\]
and define its coefficients \(\beta_{j,\alpha}(\lambda)\) by the unique expansion
\[
P_{j,\lambda}(u)-f_\lambda(c_j^0)
=
\sum_{\alpha\in\{0,\ldots,D\}^4}
\beta_{j,\alpha}(\lambda)
\prod_{\zeta\in\mathcal Z}
\left(\frac{u_\zeta-c^0_{\zeta j}}{R_{\zeta j}}\right)^{\alpha_\zeta}.
\]
For integers \(h\geq0\), define
\[
U_h(N;z)
=
\sum_{t=0}^h {h\choose t}(-z)^{h-t}\frac{(N)_t}{m^t},
\qquad
(N)_t=N(N-1)\cdots(N-t+1),
\qquad
(N)_0=1.
\]
The cell estimate, its clipped version, and the auxiliary threshold process are
\[
Z_j(\lambda)
=
f_\lambda(c_j^0)+
\sum_{\alpha\in\{0,\ldots,D\}^4}
\beta_{j,\alpha}(\lambda)
\prod_{\zeta\in\mathcal Z}
\frac{U_{\alpha_\zeta}(N_{\zeta j};c^0_{\zeta j})}
{R_{\zeta j}^{\alpha_\zeta}},
\]
\[
T_j(\lambda)
=
f_\lambda(c_j^0)+
\Pi_{[-d^{1/4}S_j,d^{1/4}S_j]}
\!\left(Z_j(\lambda)-f_\lambda(c_j^0)\right),
\qquad
\widetilde F^{\mathrm{cap}}(\lambda)
=
\sum_{j=1}^d T_j(\lambda).
\]
Here
\[
\Pi_{[r,s]}(x)=\min\{s,\max\{r,x\}\}.
\]
In this branch, average the auxiliary process conditional on the observed sample:
\[
\widehat F(\lambda)
=
E\!\left\{\widetilde F^{\mathrm{cap}}(\lambda)\mid O^{(n)}\right\}.
\]

\item In the empirical and Jackson–factorial branches, report
\[
\widehat V_b^{\mathrm{JF}}
=
\Pi_{[0,1]}
\left(
\min_{0\leq\lambda\leq1}
\{b\lambda+\widehat F(\lambda)\}
\right),
\qquad b\in I.
\]
The output uses the minimum value and requires no optimizer selection. If an implementation stores an optimizer, it stores the least element of the nonempty compact argmin set. These branch rules, the zero process on \(\{M>n\}\), and the least-minimizer convention apply under every data or treatment-effect tie.
\end{enumerate}
\end{algorithmv}

The conditional average in \cref{obj:def:jf-estimator} is an observed-sample statistic. The dual contraction transfers a uniform process bound to every budget in \(I\).

\begin{theoremv}{thm:curve-upper}[Uniform budget-curve upper bound]
Fix the overlap parameter \(\epsilon\). Suppose:
\begin{itemize}
\item \textbf{(Overlap.)} \(0<\epsilon<1/2\).
\item \textbf{(Sample and alphabet.)} \(n\geq 1\) and \(d\geq 2\).
\item \textbf{(Budget interval.)} \(b_0\in[0,1/2]\), and \(I=[b_0,1-b_0]\).
\item \textbf{(Sampling.)} For each causal law \(P\in\mathcal M_{d,\epsilon}\), the \(n\) observed records follow \cref{obj:ass:iid-sampling} for the observed margin of \(P\).
\end{itemize}
There is a constant \(C_\epsilon>0\), depending only on \(\epsilon\), such that the estimator in \cref{obj:def:jf-estimator} and the budget value in \cref{obj:def:budget-value} satisfy
\[
\sup_{P\in\mathcal M_{d,\epsilon}}
E_P\!\left[
\sup_{b\in I}
\left|\widehat V_b^{\mathrm{JF}}-V_b(P)\right|^2
\right]
\leq
C_\epsilon\min\left\{1,\frac{d}{n\log(ed)}\right\}.
\]
\end{theoremv}

The estimator controls the threshold process over the full interval of shadow prices before taking the dual minimum. Since the minimum value is Lipschitz in the supremum norm of that process, one process bound controls every budget in \(I\), including zero and full capacity when \(b_0=0\).

\subsection{Process estimates}

The upper-bound analysis begins with regularity of the cell contribution in the dual representation of \cref{obj:def:dual-process}. Bounded variation controls its dependence on the full interval of shadow prices.

\begin{lemmav}{lem:bv-lipschitz}[Bounded variation Lipschitz bound]
For \(0<\epsilon<1/2\), there is a finite constant \(C_\epsilon\ge 0\) such that, for every \(u,v\in\mathbb R_+^4\),
\[
\|f_{\cdot}(u)-f_{\cdot}(v)\|_{BV([0,1])}
\le C_\epsilon\|u-v\|_1,
\]
where \(\|h\|_{BV([0,1])}=|h(0)|+\operatorname{TV}_{[0,1]}(h)\).
\end{lemmav}

The bound in \cref{obj:lem:bv-lipschitz} supplies a common regularity measure across shadow prices. The process criterion below collects the corresponding approximation, bias, and fluctuation bounds for the estimator in \cref{obj:def:jf-estimator}.

\begin{definitionv}{def:process-handle}[Uniform process handle $\mathsf H_\epsilon$]
For fixed \(0<\epsilon<1/2\), the uniform process handle \(\mathsf H_\epsilon\) is the set of constants \(C>0\) satisfying the five groups of bounds below for the single pilot-local Jackson polynomial and clipping rule of \(\widehat V^{\mathrm{JF}}\). Membership requires the bounds simultaneously for every \(n,d\) in the Jackson–factorial branch, every observed categorical law \(P\) on the known \(d\)-cell alphabet, every pilot realization, every cell \(j\), and every \(\lambda\in[0,1]\).

Set
\[
L=\log(ed),\qquad
m=\frac n8,\qquad
K=\max\left\{2,\left\lfloor\frac L{512}\right\rfloor\right\},
\qquad
S_j=\sum_{\zeta\in\mathcal Z}R_{\zeta j},
\qquad
v_j=\sqrt{\frac{p_jL}{m}}+\frac Lm.
\]
For the expectations and \(L_2\) norms below, use the uncapped marked-Poisson experiment in which all pilot and evaluation counts \(N'_{\zeta j}\) and \(N_{\zeta j}\) are mutually independent with distribution \(\operatorname{Pois}(m q_{\zeta j})\). Define
\[
G_j
=
\bigcap_{\zeta\in\mathcal Z}
\left\{
|c_{\zeta j}-q_{\zeta j}|
\leq \frac{h_{\zeta j}}4
\right\},
\qquad
\|h\|_{BV([0,1])}
=
|h(0)|+\operatorname{TV}_{[0,1]}(h).
\]
Evaluate \(T_j\), \(\beta_{j,\alpha}\), and \(R_{\zeta j}\) in this experiment. The membership bounds are:
\begin{itemize}
\item \textbf{Coefficient envelope and exponential bound.}
\[
\sum_\alpha\|\beta_{j,\alpha}\|_{BV([0,1])}
\leq C S_j e^{12K},
\qquad
S_j^2 e^{24K+16K^2/L}
\leq C d^{1/16}S_j^2.
\]

\item \textbf{Centered process on the pilot event.}
\[
E\left\|
\sum_{j=1}^d
\left[
\mathbf1_{G_j}\{T_j-f_\cdot(q_j)\}
-
E\mathbf1_{G_j}\{T_j-f_\cdot(q_j)\}
\right]
\right\|_\infty^2
\leq C\frac d{mL}.
\]

\item \textbf{Cellwise bias.}
\[
\sup_{0\leq\lambda\leq1}
\left|ET_j(\lambda)-f_\lambda(q_j)\right|
\leq
C\left\{\sqrt{\frac{p_j}{mL}}+\frac1{mL}\right\}.
\]

\item \textbf{Cellwise envelope outside the pilot event.}
\[
\left\|
\mathbf1_{G_j^c}
\sup_{0\leq\lambda\leq1}
|T_j(\lambda)-f_\lambda(q_j)|
\right\|_{L_2}
\leq
C d^{1/4}e^{-10L}v_j.
\]

\item \textbf{Centered process outside the pilot event.}
\[
E\left\|
\sum_{j=1}^d
\left[
\mathbf1_{G_j^c}\{T_j-f_\cdot(q_j)\}
-
E\mathbf1_{G_j^c}\{T_j-f_\cdot(q_j)\}
\right]
\right\|_\infty^2
\leq C\frac d{mL}.
\]
\end{itemize}
\end{definitionv}

The criterion in \cref{obj:def:process-handle} treats pilot localization and its complement within the same process bound. A second-moment bound for the pilot radii supplies an integrability input for that criterion.

\begin{lemmav}{lem:pilot-radius-second-moment}[Second moment of pilot radii]
Let \(n\geq 1\), \(d\geq 16\), and \(m=n/8\). Let \(P\) be a categorical law on the \(d\)-cell alphabet, with category-cell masses \(q_{\zeta j}\) and cell masses \(p_j=\sum_{\zeta\in\mathcal Z}q_{\zeta j}\), and fix a cell \(j\). In the ideal Poisson experiment, the pilot counts \(N'_{\zeta j}\) are independent with means \(m q_{\zeta j}\). Let \(S_j=\sum_{\zeta\in\mathcal Z}R_{\zeta j}\) be the sum of the four pilot radii in \cref{obj:def:jf-estimator}, and set
\[
L=\log(ed),\qquad H_0=4096,\qquad
v_j=\sqrt{\frac{p_jL}{m}}+\frac{L}{m}.
\]
Then \(S_j^2\) is integrable under the ideal Poisson law, and
\[
\mathbb E[S_j^2]\leq 32H_0^2v_j^2.
\]
\end{lemmav}

The bound in \cref{obj:lem:pilot-radius-second-moment} controls the random localization scale by cell mass and the logarithmic alphabet factor. On a localization rectangle, the Jackson approximation has the following coordinatewise bound, including at its boundary.

\begin{lemmav}{lem:jackson-boundary-adaptive}[Boundary-adaptive Jackson approximation]
Let \(P_\lambda\) be the local polynomial for \(f_\lambda\) of \cref{obj:def:dual-process}, formed by the normalized order-four Jackson construction in \cref{obj:def:jf-estimator} with pilot center and radii replaced by \(c^0\) and \(R\), and with the value \(f_\lambda(c^0)\) added to the normalized polynomial evaluation. Suppose that:
\begin{itemize}
\item \textbf{(Parameters.)} The overlap parameter satisfies \(\epsilon>0\), the shadow price satisfies \(\lambda\in[0,1]\), and \(K\) is a positive integer.
\item \textbf{(Rectangle.)} For every \(\zeta\in\mathcal Z\), \(R_\zeta>0\) and \(c^0_\zeta-R_\zeta\geq0\).
\item \textbf{(Evaluation point.)} For every \(\zeta\in\mathcal Z\), \(|y_\zeta-c^0_\zeta|\leq R_\zeta\).
\end{itemize}
Then
\[
\bigl|P_\lambda(y)-f_\lambda(y)\bigr|
\leq
32\bigl\{3(1+\epsilon^{-1})+1\bigr\}
\sum_{\zeta\in\mathcal Z}
\left\{
\frac{\sqrt{(y_\zeta-c^0_\zeta+R_\zeta)(c^0_\zeta+R_\zeta-y_\zeta)}}{K}
+\frac{R_\zeta}{K^2}
\right\}.
\]
\end{lemmav}

The endpoint factors in \cref{obj:lem:jackson-boundary-adaptive} make the approximation bound sensitive to the evaluation point’s position within the pilot rectangle. To control observations outside the good-pilot event, the next lemma bounds moments of the pilot discrepancy and width together.

\begin{lemmav}{lem:pilot-bad-score-moment}[Bad pilot score moment]
Let \(n,d\geq 1\), let \(P\) be an observed categorical law on the \(d\)-cell alphabet, and fix a cell \(j\). Write \(q_{\zeta j}\) for its four category-cell masses and \(p_j=\sum_{\zeta\in\mathcal Z}q_{\zeta j}\). Set \(m=n/8\) and \(L=\log(ed)\).

Suppose the four pilot counts are independent, with \(N'_{\zeta j}\sim\operatorname{Pois}(mq_{\zeta j})\) for each \(\zeta\in\mathcal Z\). Form \(c_{\zeta j}\) and \(h_{\zeta j}\) as in \cref{obj:def:jf-estimator}, and let \(G_j\) be the quarter-halfwidth pilot event of \cref{obj:def:process-handle}. Define
\[
B_j=\sum_{\zeta\in\mathcal Z}\bigl(|c_{\zeta j}-q_{\zeta j}|+h_{\zeta j}\bigr),
\qquad
v_j=\sqrt{\frac{p_jL}{m}}+\frac{L}{m}.
\]
For each \(t\in\{1,2,4\}\),
\[
\mathbb E\!\left[\mathbf 1_{G_j^c}B_j^{\,t}\right]
\leq
4^t\,10^{80(t+1)}e^{-20L}v_j^{\,t}.
\]
\end{lemmav}

The exponential factor in \cref{obj:lem:pilot-bad-score-moment} controls the contribution of bad pilot localizations. Across cells, bounded variation gives a maximal inequality for the centered paths on either part of the pilot split.

\begin{lemmav}{lem:bv-maximal-cell-paths}[Maximal bound for centered cell paths]
Let \(n,d\) be nonnegative integers, let \(\epsilon\in\mathbb R\), and let \(P\) be a discrete law on the \(d\)-cell alphabet with category-cell masses \(q_{\zeta j}\). In the ideal count experiment of \cref{obj:def:process-handle}, the pilot and evaluation counts \(N'_{\zeta j}\) and \(N_{\zeta j}\) are mutually independent with \(\operatorname{Pois}(mq_{\zeta j})\) laws, where \(m=n/8\). Choose the pilot events uniformly across cells: either \(E_j=G_j\) for every \(j\), or \(E_j=G_j^c\) for every \(j\), with \(G_j\) as in \cref{obj:def:process-handle}. For each cell \(j\), define the path on \([0,1]\) by
\[
W_j(\lambda)=\mathbf 1_{E_j}\bigl\{T_j(\lambda)-f_\lambda(q_j)\bigr\},
\]
where \(T_j\) is the clipped cell estimate of \cref{obj:def:jf-estimator} and \(f_\lambda\) is defined in \cref{obj:def:dual-process}. Assume:

\begin{itemize}
\item \textbf{(Regularity.)} For every count table and cell \(j\), \(W_j\) is continuous in \(\lambda\). For each \(j\), the count-table map into \(C([0,1])\) is measurable and Bochner integrable.
\item \textbf{(Finite variation.)} For each \(j\), \(\operatorname{TV}_{[0,1]}(W_j)<\infty\) almost surely.
\item \textbf{(Second moments.)} For each \(j\), \(\bigl(\|W_j\|_\infty+\operatorname{TV}_{[0,1]}(W_j)\bigr)^2\) is integrable.
\item \textbf{(Independence.)} The paths \(W_1,\ldots,W_d\) are independent.
\end{itemize}

Then
\[
\mathbb E\left\|\sum_{j=1}^{d}\bigl(W_j-\mathbb EW_j\bigr)\right\|_\infty^2
\leq 16384\sum_{j=1}^{d}
\mathbb E\left[\bigl(\|W_j\|_\infty+\operatorname{TV}_{[0,1]}(W_j)\bigr)^2\right].
\]
\end{lemmav}

By \cref{obj:lem:bv-maximal-cell-paths}, cellwise path moments yield control uniform in the shadow price. The reported estimator averages over auxiliary draws from an observed sample; the following identity relates that average to the uncapped count experiment.

\begin{lemmav}{lem:capped-marked-count-identity}[Capped marked-count integral identity]
Fix \(n,d\in\mathbb N\) and an observed categorical law \(P\) on the \(d\)-cell alphabet, with category-cell masses \(q_{\zeta j}\). Let \(O^{(n)}=(O_1,\ldots,O_n)\sim P^{\otimes n}\). For \(M\in\{0,\ldots,n\}\), a permutation \(\sigma\) of the \(n\) record indices, and pilot/evaluation marks \(\omega\in\{\mathrm{pilot},\mathrm{evaluation}\}^n\), write \((N',N)(O^{(n)},\sigma,M,\omega)\) for the pilot and evaluation count tables formed from the first \(M\) permuted records according to \cref{obj:def:jf-estimator}. Define
\[
w_n(M)=\frac{\Pr\{\operatorname{Pois}(n/4)=M\}}{n!\,2^n}.
\]
In the uncapped marked-Poisson experiment, \(M\sim\operatorname{Pois}(n/4)\); conditional on \(M\), draw \(M\) independent records from \(P\) and mark each independently as pilot or evaluation with probability \(1/2\), yielding count tables \((N',N)\).

For every real-valued function \(g\) of pairs of count tables that is integrable under the ideal law of \cref{obj:def:process-handle}, whose pilot and evaluation category-cell counts are independent \(\operatorname{Pois}((n/8)q_{\zeta j})\) variables,
\[
\mathbb E_{O^{(n)}\sim P^{\otimes n}}\left[
\sum_{M=0}^{n}\sum_{\sigma}\sum_{\omega\in\{\mathrm{pilot},\mathrm{evaluation}\}^n}
w_n(M)\,
g\bigl((N',N)(O^{(n)},\sigma,M,\omega)\bigr)
\right]
=
\mathbb E_{\mathrm{uncapped}}\left[
\mathbf 1\{M\leq n\}\,g(N',N)
\right].
\]
\end{lemmav}

The identity in \cref{obj:lem:capped-marked-count-identity} supplies the count-law link for the conditional average. Together with the radius, approximation, bad-pilot, and maximal-path bounds, it supports a common constant for the process criterion and the threshold-process bounds.

\begin{theoremv}{thm:integrated-process-certificate}[Uniform integrated process bounds]
Fix an overlap parameter \(0<\epsilon<1/2\). There is a single constant \(C_\epsilon>0\), belonging to the uniform process handle in \cref{obj:def:process-handle}, for which the following conclusions hold whenever:

\begin{itemize}
\item \textbf{(Sample and alphabet.)} \(n\geq1\), \(d\geq16\), and \(d/\log(ed)\leq n\).
\item \textbf{(Sampling.)} The law \(\mu\) of the \(n\) observed records satisfies \cref{obj:ass:iid-sampling} for the observed margin of \(P\).
\item \textbf{(Causal law.)} The full-data law \(P\) belongs to \(\mathcal M_{d,\epsilon}\) of \cref{obj:def:model-class}.
\end{itemize}

Write \(L=\log(ed)\), \(m=n/8\), and \(v_j=\sqrt{p_jL/m}+L/m\). Simultaneously for every cell \(j\) and every \(\lambda\in[0,1]\), the absolute bias of the uncapped Jackson cell estimate of \(f_\lambda(q_j)\) is at most
\[
C_\epsilon\left\{\sqrt{\frac{p_j}{nL}}+\frac{1}{nL}\right\}.
\]
The expected squared supremum of the centered good-pilot error sum is at most \(C_\epsilon d/(nL)\). The uncentered bad-pilot cell envelope satisfies
\[
\left\|\mathbf 1_{G_j^c}\sup_{\lambda\in[0,1]}
\left|T_j(\lambda)-f_\lambda(q_j)\right|\right\|_{L_2}
\leq C_\epsilon d^{1/4}e^{-10L}v_j,
\]
and the expected squared supremum of the centered bad-pilot error sum is at most \(C_\epsilon d/(mL)\). Finally, the ideal and averaged versions of the dual threshold process in \cref{obj:def:jf-estimator} each have expected squared supremum error relative to \(F_P\) at most \(C_\epsilon d/(nL)\). The same \(C_\epsilon\) applies to all these bounds.
\end{theoremv}

The final process bounds in \cref{obj:thm:integrated-process-certificate} apply to an entire interval of shadow prices. Through the dual representation in \cref{obj:prop:dual-representation}, they provide the simultaneous control used in \cref{obj:thm:curve-upper}, including budgets selected after the curve is estimated. The average is formed at the threshold-process level before optimization.

Two boundedness facts establish that the empirical and conditionally averaged threshold processes have finite values throughout the shadow-price interval.

\begin{lemmav}{lem:empirical-threshold-process-bounded}[Bounded empirical threshold process]
For any \(n,d\in\mathbb N\), any real overlap parameter \(\epsilon\), and any sample \(O^{(n)}\) of \(n\) observations on the \(d\)-cell alphabet, let
\[
\widehat F(\lambda)=\sum_{j=1}^{d}f_\lambda(\widehat q_j)
\]
be the empirical threshold process of \cref{obj:def:jf-estimator}, with \(f_\lambda\) as in \cref{obj:def:dual-process}. There exists \(B\in\mathbb R\) such that
\[
|\widehat F(\lambda)|\leq B
\qquad\text{for every }\lambda\in[0,1].
\]
\end{lemmav}

Conditional averaging over the auxiliary Poisson count preserves boundedness for the threshold process used by the reported estimator.

\begin{lemmav}{lem:averaged-threshold-process-bounded}[Bounded averaged threshold process]
Fix any observed sample \(O^{(n)}\) of \(n\) records from a \(d\)-cell alphabet. Suppose:
\begin{itemize}
\item \textbf{(Overlap parameter.)} \(\epsilon>0\).
\item \textbf{(Sample size.)} \(n\geq 1\).
\item \textbf{(Alphabet size.)} \(d\geq 1\).
\end{itemize}
The conditionally averaged threshold process \(\widehat F(\lambda)\) of \cref{obj:def:jf-estimator}, whose Poisson\((n/4)\) weights are summed over \(M=0,\ldots,n\), is bounded on \([0,1]\): there exists \(B\in\mathbb R\) such that
\[
|\widehat F(\lambda)|\leq B
\qquad\text{for every }\lambda\in[0,1].
\]
\end{lemmav}

For the lower bound, a simplex records the masses assigned to pairs of covariate cells. The paired family then isolates the effect of allocating a fixed treatment capacity across unequal benefits.

\begin{definitionv}{def:paired-family}[Paired causal family $\mathcal P_k^{\mathrm{pair}}$]
The paired causal family \(\mathcal P_k^{\mathrm{pair}}\) consists of causal completions constructed as follows. For \(k\geq1\), choose \(r=(r_1,\ldots,r_k)\) in the simplex and \(\theta_i\in[-1/8,1/8]\). For each \(i\), create cells \((i,+)\) and \((i,-)\) with masses \(r_i/2\), and set
\[
e_{i,\sigma}=\frac12,\qquad
\mu_{0,i,\sigma}=\frac14,\qquad
\mu_{1,i,\sigma}=\frac12+\sigma\theta_i,
\qquad
\sigma\in\{+1,-1\}.
\]
Generate the potential outcomes independently of \(A\) conditional on \(X\).
\end{definitionv}

The simplex \(\Delta_k\) in \cref{obj:synth_10} supplies the pair-mass vector for \cref{obj:def:paired-family}.

\begin{definitionv}{synth_10}[Probability simplex on \(k\) categories]
For an integer \(k\geq1\), define
\[
\Delta_k=\left\{r=(r_1,\ldots,r_k)\in[0,1]^k:\sum_{i=1}^k r_i=1\right\}.
\]
\end{definitionv}

With the pair masses specified, the causal family fixes the propensity and control mean while varying the treatment mean within each pair. Treatment effects remain positive while their ranking depends on the contrast. The capacity and sampling identities make this allocation problem explicit.

\begin{propositionv}{prop:paired-capacity-identity}[Paired capacity and sample identity]
Let \(k\geq 1\) and \(0<\epsilon<1/2\). The following statements hold.
\begin{itemize}
\item \textbf{(Capacity identity.)} For every \(P\in\mathcal P^{\mathrm{pair}}_k\) and every pair-mass vector \(r\in\Delta_k\) and contrast vector \(\theta\in[-1/8,1/8]^k\) for which \(P\) satisfies the paired-completion conditions of \cref{obj:def:paired-family}, the law \(P\) belongs to \(\mathcal M_{2k,\epsilon}\) of \cref{obj:def:model-class}. Every occupied cell has treatment effect \(\tau_j\in[1/8,3/8]\). The half-population budget is attained by a policy \(\pi\in\Pi_{1/2}(P)\) satisfying
\[
\sum_j p_j\pi_j=\frac12,
\qquad
V_{1/2}(P)=B(P)+\sum_j p_j\pi_j\tau_j.
\]
For the budget values of \cref{obj:def:budget-value},
\[
V_{1/2}(P)
=\frac38+\frac12\sum_{i=1}^k r_i|\theta_i|
=\frac18+F_P\!\left(\frac14\right),
\qquad
V_1(P)=\frac12.
\]

\item \textbf{(Exact sample reduction.)} For every integer \(n\geq0\), there is a Markov kernel from pairs of \(n\)-record samples on the \(k\)-point alphabet to \(n\) observed records on the \(2k\)-cell alphabet with the following property. For every pair \(R,S\in\Delta_k\), the simplex of \(k\)-point laws, set
\[
r_i=\frac{R_i+S_i}{2},
\qquad
\theta_i=
\begin{cases}
0, & R_i+S_i=0,\\[2pt]
\dfrac{R_i-S_i}{8(R_i+S_i)}, & R_i+S_i>0.
\end{cases}
\]
The paired law determined by \((r,\theta)\) belongs to \(\mathcal P^{\mathrm{pair}}_k\). Applying the kernel to independent \(n\)-record samples from \(R\) and \(S\) yields exactly the law of \(n\) independent observed records from the observed margin of that paired law.
\end{itemize}
\end{propositionv}

For two sampling laws \leanref{sym:R}{\ensuremath{R}} and \leanref{sym:S}{\ensuremath{S}}, the exact sample reduction gives a half-budget value of \(3/8+\|R-S\|_1/32\), while the full-budget value is \(1/2\). Thus the half-budget comparison carries a two-sample distance problem within the causal model. The statistical lower-bound input is a squared-error bound for the two-sample distance.

The following bound states the two-sample risk at the scale used by the reduction.

\begin{lemmav}{lem:two-sample-l1-lower}[Two-sample \(L_1\) lower bound]*
There is a universal constant \(c>0\) such that, for every integer \(k\geq 2\) and \(1\leq n\leq k^2\), the minimax squared-error risk for estimating \(\|R-S\|_1\) from \(n\) independent observations from each of \(R,S\in\Delta_k\), the probability simplex on \(k\) categories, satisfies
\[
\inf_{\widehat L}\sup_{R,S\in\Delta_k}
E_{R,S}\!\left[(\widehat L-\|R-S\|_1)^2\right]
\geq c\min\left\{1,\frac{k}{n\log(ek)}\right\}.
\]
\end{lemmav}
% CAUSALSMITH-CITED-SCOPE-BEGIN lem:two-sample-l1-lower
\verificationfootnotetext{\textbf{Formalization scope.} This result uses the published conclusion \citep{JiaoHanWeissman2018} (Theorem 3 and the unknown-both-distributions reduction following Theorem 6). The cited source proof is not formalized here: Lean verifies its use and all remaining steps. If that cited conclusion is false, the portions of this result that depend on it are not certified.}
% CAUSALSMITH-CITED-SCOPE-END lem:two-sample-l1-lower

The published intermediate-regime lower-bound input for \cref{obj:lem:two-sample-l1-lower} is given by \citet[Theorem 3 and the unknown-both-distributions reduction following Theorem 6]{JiaoHanWeissman2018}. Combined with the exact sample reduction in \cref{obj:prop:paired-capacity-identity}, the lemma supplies the half-budget lower bound in \cref{obj:thm:capacity-active-lower}. That bound and the upper bound in \cref{obj:thm:curve-upper} yield the matched risk statement in \cref{obj:thm:matched-curve-frontier} under its stated conditions.

The theorem statements and internal proofs are machine-checked in Lean 4. The sampling and causal conditions in \cref{obj:ass:iid-sampling,obj:ass:consistency,obj:ass:conditional-exchangeability,obj:ass:overlap} are model assumptions. The published conclusion cited for \cref{obj:lem:two-sample-l1-lower} is a mathematical input from the literature.
